% Auto-generated by scripts/make_perfacet_tables.py
\begin{table}[t]
\centering
\small
\setlength{\tabcolsep}{4pt}
\begin{tabular}{@{}l rrrr l@{}}
\toprule
\textbf{Concept (type)} & \textbf{KW} & \textbf{+Gem} & \textbf{v1} & \textbf{v2} & \textbf{Synonyms folded in (stage)} \\
\midrule
husband (Abbreviation) & 442 & 474 & 480 & 480 & Gemma: dh, hubby, hubs, hubbie; v1: hubby, dh; v2: hubs \\
tangy (Synonym cluster) & 41 & 32 & 61 & 73 & Gemma: tart, tang; v1: tang, tart; v2: tartness, zesty, zingy \\
\bottomrule
\end{tabular}
\caption{Two illustrative facets across the pipeline (recipes of 500). \emph{husband} is an abbreviation cluster: Gemma resolves most DH/hubby internally, v1 folds the residual. \emph{tangy} is a semantic cluster: Gemma is selective (the bare word can be a negation or modify a condiment), v1 folds tang/tart, and v2 adds the long tail tartness/zesty/zingy. Counts are recipe unions, so a frequent variant that co-occurs with the canonical adds few new recipes; the value is one complete facet.}
\label{tab:perfacet-example}
\end{table}
